% !TeX root=../main.tex

\chapter{نمونه های اجرا}
در این قسمت یک نمونه پرسش که هم از مدل پایه (آیا اکسپنس) هم از مدل های معرفی شده پرسیده شده آمده است، لازم به ذکر است که تمامی سوالات پرسیده شده از مدل ها به طور کامل در مخازن گیت هاب به صورت فایل اکسل موجود است و اینجا تنها به چند نمونه از آنها اشاره شده است؛ در نظر داشته باشید همانطور که گفته شد اگر از مدل ها درخواست زنجیره افکار شود ممکن است نتایج هر مرتبه متفاوت باشد در نمونه های زیر به نوشتن نتایج یکی از اجراها اکتفا شده است.
\footnote{
آدرس مخازن گیت هاب همانطور که پیشتر اشاره شد
\lr{github.com/mehrdadghassabi/gaokerena-V}
و
\lr{github.com/mehrdadghassabi/gaokerena-R}
میباشند.
}
\section{
نمونه های اجرای گائوکرنا وی در مقایسه با مدل پایه
}
\textbf{
درخواست
\footnote{\lr{prompt}}
:
}

به پرسش زیر به پاسخ بدهید و تنها از زبان فارسی استفاده کنید!

\#\# پرسش

کدام یک از گزینه های زیر در مورد عملکرد انسولین درست است؟ 

\#\# گزینه ها

\lr{[A] }
:
گليکوژنوليز(تجزیه گليکوژن) در سلول کبدی را افزایش می دهد

\lr{[B] }
:
بتا اکسيداسيون اسيدهای چرب را در ميتوکندریسلول های کبدی افزایش می دهد

\lr{[C] }
:
غلظت پلاسمایی کلسترول و فسفوليپيدها را افزایش می دهد

\lr{[D] }
:
آنزیمليپو پروتئين ليپاز در دیواره مویرگ بافت چربی را فعال می کند

\textbf{
مطرح شده در:
}

پرسش هفدهم کنکور علوم پایه پزشکی شهریور
\lr{1402}

\textbf{
گزینه صحیح:
}

گزینه
\lr{[D]}

\textbf{
پاسخ آیا اکسپنس:
}

پاسخ صحیح **نیست هیچ‌کدام از گزینه‌ها**.

انسولین هورمونی است که توسط لوزالمعده ترشح می‌شود و نقش اصلی آن در تنظیم متابولیسم گلوکز است. عملکرد اصلی آن **کاهش غلظت گلوکز خون** از طریق دو مکانیسم اصلی است:

\lr{1}
.
**افزایش جذب گلوکز توسط سلول‌ها:** انسولین به سلول‌ها اجازه می‌دهد تا گلوکز را از خون به عنوان منبع انرژی جذب کنند.

\lr{2}
.
 **مهار گلیکوجنولیز:** انسولین تولید گلیکوجن (فرم ذخیره گلوکز) در کبد و عضلات را مهار می‌کند.

گزینه‌های ارائه شده با عملکرد واقعی انسولین مطابقت ندارند:

\lr{* [A]}
غلظت گلیکوجن را افزایش می‌دهد، نه کاهش.

\lr{* [B]}
بتا اکسیداسیون اسیدهای چرب را افزایش می‌دهد، که ارتباط مستقیمی با انسولین ندارد.

\lr{* [C]}
انسولین غلظت کلسترول و فسفولیپید را **کاهش** می‌دهد، نه افزایش.

\lr{* [D]}
انسولین لیپاز را فعال نمی‌کند، بلکه متابولیسم چربی را به طور کلی تنظیم می‌کند.

\textbf{
پاسخ گائوکرنا وی:
}

پاسخ صحیح گزینه
\lr{[D]}
است.

انسولین هورمونی است که توسط لوزالمعده تولید می‌شود و نقش مهمی در تنظیم قند خون دارد. عملکرد اصلی انسولین، افزایش جذب گلوکز توسط سلول‌ها است. این هورمون به طور مستقیم بر روی سلول‌های هدف مانند سلول‌های کبد، عضلات و بافت چربی تأثیر می‌گذارد.

گزینه
\lr{[D]}
که «آنزیم لیپو پروتئین لیپاز در دیواره مویرگ بافت چربی را فعال می‌کند» درست است. انسولین باعث افزایش فعالیت آنزیم لیپو پروتئین لیپاز
\lr{(LPL)}
در دیواره مویرگ بافت چربی می‌شود. این آنزیم، چربی‌های ذخیره شده در بافت چربی را تجزیه کرده و باعث آزاد شدن اسیدهای چرب به جریان خون می‌شود. این فرآیند به عنوان لیپولیز شناخته می‌شود و به بدن کمک می‌کند تا از چربی‌ها به عنوان منبع انرژی استفاده کند.

گزینه
\lr{[A]}
غلط است. انسولین باعث کاهش تجزیه گلوکز (گلیکوجنولیز) در کبد نمی‌شود.

گزینه
\lr{[B]}
غلط است. انسولین باعث افزایش بتا اکسیداسیون اسیدهای چرب در میتوکندری سلول‌های کبدی نمی‌شود.

گزینه
\lr{[C]}
غلط است. انسولین باعث افزایش غلظت پلاسمایی کلسترول و فسفولیپیدها نمی‌شود.

بنابراین، گزینه
\lr{[D]}
تنها گزینه صحیح در مورد عملکرد انسولین است.







\section{
نمونه های اجرای گائوکرنا آر در مقایسه با مدل پایه
}

\textbf{
درخواست
\footnote{\lr{prompt}}
:
}

به پرسش زیر قدم به قدم فکر کنید و زنجیره افکار 
\lr{(chain of thought)}
خود برای رسیدن به پاسخ را به طور کامل شرح دهید و تنها یک گزینه را انتخاب کنید.

\#\# پرسش:

بیماری با آسیب وسیع درد و طرف نواحی داخلی قشر پس سری و گیجگاهی کدام علایم را دارد؟ 


\#\# گزینه ها

[A] : هيپوتونی در سمت چپ

[B] : سفتی در انگشتان هر دو دست

[C] : غفلت درنيمه چپ بدن

[D] : عدم تشخيص چهره ها 

\textbf{
مطرح شده در:
}

پرسش ششم کنکور علوم پایه پزشکی شهریور
\lr{1402}

\textbf{
گزینه صحیح:
}

گزینه
\lr{[D]}


\textbf{
پاسخ آیا اکسپنس:
}

**زنجیره افکار:** 

برای پاسخ به این پرسش، ما باید درک کنیم که نواحی قشر پس سری و گیجگاهی کجاست و آسیب به آن چه اثراتی می‌تواند داشته باشد.  

- **قشر پس سری و گیجگاهی:** این نواحی از قشر مغز مسئول عملکردهای مهمی مانند درک حس،  کنترل حرکات، ادراک فضایی، حافظه  و زبان است.

- **علائم احتمالی آسیب:** با توجه به توصیف "آسیب وسیع درد و آسیب در نواحی داخلی قشر پس سری و گیجگاهی"،  می‌توانیم انتظار علائم زیر را داشت:

    - **غفلت در نیمکره چپ:**  آسیب به قشر مغز می‌تواند منجر به  اختلال در پردازش اطلاعات حسی و حرکتی در نیمکره چپ شود که  علائمی مانند غفلت یا بی‌توجهی در آن منطقه ایجاد کند.

    - **عدم تشخیص چهره‌ها:**   ناحیه‌ای از قشر به نام "قلمروی چهره" در ناحیه گیجگاهی واقع شده است. آسیب به این ناحیه می‌تواند منجر به مشکلات در شناسایی چهره‌ها شود.

**انتخاب گزینه:**

بنابراین، با توجه به زنجیره افکار فوق، پاسخ صحیح
\lr{**[C]}
غفلت در نیمکره چپ بدن** است. 


* **دلیل رد گزینه‌های دیگر:**

    \lr{* **[A]}
    هيپوتونی در سمت چپ:** هيپوتونی ضعف عضلانی است که بیشتر با آسیب به سیستم عصبی حرکتی مرتبط است، نه به طور مستقیم نواحی قشر پس سری و گیجگاهی.
    
    \lr{* **[B]}
    سفتی در انگشتان هر دو دست:**  این علامت معمولا با اختلالات سیستم عصبی مرتبط است، نه به طور خاص آسیب قشر پس سری و گیجگاهی.
    
    \lr{* **[D]}
    عدم تشخیص چهره ها:** این علائم مستقیماً با منطقه آسیب  دیده مطابقت دارد.

\textbf{
پاسخ گائوکرنا آر:
}

برای پاسخ به این پرسش، باید به ارتباط بین نواحی آسیب‌دیده در قشر مغز و علایمی که هر گزینه توصیف می‌کند، توجه کنیم. بیماری با آسیب درد وسیع در نواحی داخلی قشر پس سری و گیجگاهی، احتمالاً با علائمی مرتبط با عملکرد این نواحی همراه است.

**قشر پس سری 
\lr{(Posterior Parietal Cortex)}
:**

- مسئولیت پردازش اطلاعات حسی، به ویژه از طریق بینایی و فاصله را بر عهده دارد.

- آسیب به این ناحیه می‌تواند منجر به مشکلاتی در تشخیص اشکال، جهت‌یابی، و درک فضا شود.

**قشر گیجگاهی 
\lr{(Temporal Lobes)}
:**

- در پردازش زبان، حافظه، و تشخیص چهره نقش دارد.

- آسیب‌های این ناحیه می‌تواند منجر به اختلالات زبانی، مشکلات حافظه‌ای و شناختی، و حتی ناتوانی در تشخیص چهره‌ها شود.

با توجه به این اطلاعات، گزینه‌های را یکی پس از دیگری بررسی می‌کنیم:

\lr{[A]}
**هیپوتونی در سمت چپ:** هیپوتونی به معنای ضعف عضلانی است و معمولاً با آسیب‌های موتوریک مرتبط است. این علائم با نواحی آسیب‌دیده در پرسش مرتبط نیست.

\lr{[B]}
**سفتی در انگشتان هر دو دست:** این علائم ممکن است با آسیب‌های موتورik مرتبط باشد، اما با نواحی ذکر شده در پرسش مطابقت ندارد.

\lr{[C]}
**غفلت در نیم چپ بدن:** غفلت
\lr{(Neglect)}
یک حالت حسی یا حرکتی است که معمولاً در اثر آسیب به نیمکره چپ مغز رخ می‌دهد. این پدیده شامل بی‌توجهی به محیط، اشیاء، یا قسمت‌های بدن در یک طرف است. با توجه به اینکه پرسش از آسیب به نواحی داخلی قشر پس سری و گیجگاهی سخن می‌گوید، این گزینه ممکن است مرتبط باشد، به ویژه اگر آسیب به سمت چپ مغز محدود باشد.

\lr{[D]}
**عدم تشخیص چهره‌ها:** این علائم مستقیماً با نواحی آسیب‌دیده در پرسش ارتباط دارد. قشر گیجگاهی در تشخیص چهره‌ها نقش دارد، بنابراین آسیب به این ناحیه می‌تواند منجر به ناتوانی در تشخیص چهره‌ها شود.

بنابراین، با توجه به زنجیره افکار فوق، **گزینه
\lr{[D]}
عدم تشخیص چهره‌ها** پاسخ صحیح است.

